\documentclass[conference]{IEEEtran}
\IEEEoverridecommandlockouts
% The preceding line is only needed to identify funding in the first footnote. If that is unneeded, please comment it out.
\usepackage{cite}
\usepackage{amsmath,amssymb,amsfonts}
\usepackage{algorithmic}
\usepackage{graphicx}
\usepackage{textcomp}
\usepackage{xcolor}
\usepackage{booktabs}
\usepackage{url}
\usepackage{microtype}

\def\BibTeX{{\rm B\kern-.05em{\sc i\kern-.025em b}\kern-.08em
    T\kern-.1667em\lower.7ex\hbox{E}\kern-.125emX}}

\begin{document}

\title{A Leak-Free Machine Learning Framework for Rapid Detection and Diagnostic Hierarchy of Milk Adulteration}

\author{\IEEEauthorblockN{Asmit Das}
\IEEEauthorblockA{\textit{Dept. of Computer Science \& Eng.} \\
\textit{Manipal Institute of Technology (MIT) Bengaluru}\\
\textit{Manipal Academy of Higher Education}\\
Bengaluru, India \\
asmit.das@learner.manipal.edu}
\and
\IEEEauthorblockN{Suprathik Narendra}
\IEEEauthorblockA{\textit{Dept. of Computer Science \& Eng.} \\
\textit{Manipal Institute of Technology (MIT) Bengaluru}\\
\textit{Manipal Academy of Higher Education}\\
Bengaluru, India \\
suprathik.narendra@learner.manipal.edu}
\and
\IEEEauthorblockN{Nadia Cecil}
\IEEEauthorblockA{\textit{Dept. of Computer Science \& Eng.} \\
\textit{Manipal Institute of Technology (MIT) Bengaluru}\\
\textit{Manipal Academy of Higher Education}\\
Bengaluru, India \\
nadia.cecil@learner.manipal.edu}
\and
\IEEEauthorblockN{Dr. S. Vaishnavi}
\IEEEauthorblockA{\textit{Dept. of Computer Science \& Eng.} \\
\textit{Manipal Institute of Technology (MIT) Bengaluru}\\
\textit{Manipal Academy of Higher Education}\\
Bengaluru, India \\
vaishnavi.s@manipal.edu}
}

\maketitle

\begin{abstract}
Food adulteration is a critical global challenge impacting nutritional security, public health, and fair trade. In developing agrarian economies, conventional analytical screening protocols such as High-Performance Liquid Chromatography (HPLC) and Fourier-Transform Infrared (FTIR) spectroscopy are economically and logistically unviable at rural collection centers. In this paper, we propose an end-to-end, leak-free machine learning framework for rapid, reagent-free screening of milk adulteration using seven routine physicochemical parameters (pH, temperature, taste, odor, fat, turbidity, and color). Employing strict Test-Driven Development (TDD), we eliminate critical statistical flaws prevalent in conventional pipelines by enforcing stratified dataset partitioning prior to any preprocessing transformations. Within each fold, missing values are median-imputed, extreme leverage points are bounded via adaptive Interquartile Range (IQR) Winsorization, and features are standardized. We benchmark five optimized classifiers against a dummy baseline using 5-Fold Stratified Cross-Validation on 847 training records. On a locked, held-out test cohort of 212 samples evaluated exactly once, our tuned Support Vector Machine (RBF kernel, $C=100.0, \gamma=\text{'scale'}$) achieves an accuracy of \textbf{99.53\%}, an $F_1$ score of \textbf{0.9969}, an ROC-AUC of \textbf{1.0000}, and a \textbf{100.00\% Recall} with \textbf{zero false negatives}. Model-agnostic permutation importance and white-box decision tree induction establish pH ($\Delta F_1 = 0.2088$) and turbidity ($\Delta F_1 = 0.0839$) as the primary diagnostic indicators. Extensive ablation studies validate that IQR Winsorization preserves critical adulterant signals that row-deletion techniques discard.
\end{abstract}

\begin{IEEEkeywords}
Milk Adulteration, Food Safety, Machine Learning, Data Leakage Prevention, Support Vector Machines, Feature Importance, Interpretability, Ablation Analysis.
\end{IEEEkeywords}

\section{Introduction}
Milk is an indispensable dietary component consumed by billions worldwide. However, due to fragmented rural supply chains and commercial incentives, milk remains one of the most widely counterfeited food commodities globally~\cite{fssai2019}. Common adulteration vectors include volumetric dilution with contaminated water, the addition of hazardous chemical neutralizers (e.g., sodium hydroxide, sodium bicarbonate) to suppress lactic acid souring during non-refrigerated transit, and the addition of artificial lipid emulsions (detergents, starch, urea) to mimic fat content and density.

While advanced analytical chemistry laboratories utilize mass spectrometry and chromatographic assays for regulatory validation, these techniques require capital investments exceeding \$50,000, specialized cleanrooms, and multi-day turnaround times~\cite{poonia2017}. Conversely, handheld strip tests and lactometers are notoriously susceptible to compensatory counter-adulteration~\cite{sharma2020}.

Machine learning (ML) models operating on multi-parametric sensor data offer a compelling compromise: rapid, non-destructive, low-cost screening~\cite{saadatmand2022}. Nonetheless, a substantial portion of published food-adulteration ML literature suffers from fundamental methodological defects:
\begin{enumerate}
    \item \textbf{Pervasive Data Leakage}: Standardizing, imputing, or pruning outliers across the entire corpus prior to train/test partitioning, artificially inflating reported accuracies~\cite{kaufman2012}.
    \item \textbf{Harmful Outlier Deletion}: Treating severe physical deviations as measurement noise, inadvertently removing the very contaminated samples the model is tasked with identifying.
    \item \textbf{Black-Box Opacity}: Presenting high-dimensional models without explaining which physical parameters govern classifications.
\end{enumerate}

In this work, we present an end-to-end, test-driven machine learning framework designed to systematically overcome these deficiencies.

\section{Related Work}
Recent investigations have applied machine learning to food quality screening. Sengupta et al.~\cite{sengupta2016} explored artificial neural networks for adulteration detection using optical spectroscopy, achieving high precision but requiring expensive spectrophotometric hardware. Kar et al.~\cite{kar2018} examined electrical impedance spectroscopy for detecting water and detergent in bovine milk. In tabular physicochemical screening, previous studies frequently applied decision trees and Random Forests directly to raw datasets~\cite{kumar2021}. 

However, as highlighted by Kaufman et al.~\cite{kaufman2012}, the majority of published tabular ML studies fail to encapsulate preprocessing inside cross-validation folds, causing subtle data leakage. Furthermore, prior studies rarely perform formal ablation analysis to justify their preprocessing transformations or quantify false negative rates in clinical and safety contexts.

\section{System Architecture \& Methodology}

\begin{figure}[htbp]
\centering
\fbox{
\parbox{0.92\linewidth}{
\centering\small
\textbf{Raw Dataset ($N=1,059$ Samples)} \\
$\downarrow$ \\
\textbf{Stratified Split (80/20, seed=42)} \\
\textit{Training Set: 847 samples} $\mid$ \textit{Locked Test Set: 212 samples} \\
$\downarrow$ \\
\textbf{Scikit-Learn Pipeline Architecture (Within Folds)} \\
$\left[\begin{array}{l}
\text{1. Median Imputation} \\
\text{2. Adaptive IQR Clipper (Winsorization)} \\
\text{3. StandardScaler} \\
\text{4. Supervised Classifier (SVM / RF / GB / k-NN / DT)}
\end{array}\right]$ \\
$\downarrow$ \\
\textbf{5-Fold Stratified Cross-Validation Tuning} \\
$\downarrow$ \\
\textbf{One-Shot Final Evaluation on Locked Test Cohort} \\
\textit{Accuracy: 99.53\% $\mid$ Recall: 100.00\% $\mid$ False Negatives: 0}
}
}
\caption{End-to-End Leak-Free Machine Learning Architecture.}
\label{fig:architecture}
\end{figure}

\subsection{Dataset Specification and Label Formulation}
We utilize the benchmark Milk Quality Dataset containing 1,059 physical lab measurements. Each sample is described by seven features:
\begin{equation}
\mathbf{x} = [\text{pH}, T, \text{taste}, \text{odor}, \text{fat}, \text{turbidity}, \text{color}]^T
\end{equation}
The target classification $y \in \{0, 1\}$ is formulated per international regulatory criteria (Decision D-04):
\begin{itemize}
    \item \textbf{Class 0 (Pure / Grade High)}: $N=256$ ($24.17\%$). Compliant with biological milk standards.
    \item \textbf{Class 1 (Adulterated / Substandard / Grade Low \& Medium)}: $N=803$ ($75.83\%$). Chemically neutralized, bacterially soured, or diluted samples.
\end{itemize}

\subsection{Mathematical Leakage Guard \& Preprocessing}
To mathematically guarantee zero data leakage, partitioning occurs prior to any parameter estimation:
\begin{equation}
\mathcal{D}_{\text{train}} \cap \mathcal{D}_{\text{test}} = \emptyset, \quad |\mathcal{D}_{\text{train}}| = 847, \quad |\mathcal{D}_{\text{test}}| = 212
\end{equation}

Inside each training fold $k$, transformations are strictly encapsulated within an \texttt{sklearn.pipeline.Pipeline}:
\begin{enumerate}
    \item \textbf{Median Imputation}:
    \begin{equation}
    \tilde{x}_j = \text{median}(\{x_{i,j} \mid i \in \mathcal{D}_{\text{train}}\})
    \end{equation}
    \item \textbf{Adaptive IQR Winsorization}:
    For feature $j$, let $Q_{1,j}$ and $Q_{3,j}$ denote the 25th and 75th empirical percentiles of the training fold. The interquartile range is $\text{IQR}_j = Q_{3,j} - Q_{1,j}$. The clipping bounds are defined as:
    \begin{equation}
    L_j = Q_{1,j} - 1.5 \cdot \text{IQR}_j, \quad U_j = Q_{3,j} + 1.5 \cdot \text{IQR}_j
    \end{equation}
    The transformation function $f_{\text{clip}}(x_{i,j})$ projects values into the valid hypercube:
    \begin{equation}
    f_{\text{clip}}(x_{i,j}) = \min(\max(x_{i,j}, L_j), U_j)
    \end{equation}
    \item \textbf{Z-Score Standardization}:
    \begin{equation}
    z_{i,j} = \frac{f_{\text{clip}}(x_{i,j}) - \mu_{j,\text{train}}}{\sigma_{j,\text{train}}}
    \end{equation}
\end{enumerate}

\section{Experimental Results \& Model Benchmarks}

\subsection{Baseline vs. Tuned Cross-Validation Benchmarks}
We evaluate six architectures under 5-Fold Stratified Cross-Validation on $\mathcal{D}_{\text{train}}$. Hyperparameters were tuned via \texttt{GridSearchCV} and \texttt{RandomizedSearchCV}. As shown in Table~\ref{tab:cv_benchmarks}, all supervised classifiers substantially outperform the dummy baseline.

\begin{table*}[htbp]
\caption{5-Fold Stratified Cross-Validation Benchmarks on Training Cohort ($N=847$)}
\label{tab:cv_benchmarks}
\centering
\begin{tabular}{llccccc}
\toprule
\textbf{Estimator} & \textbf{Tuning Strategy} & \textbf{Optimal Hyperparameters} & \textbf{Mean CV $F_1$} & \textbf{Mean CV Acc.} & \textbf{Mean CV Recall} & \textbf{Mean CV ROC-AUC} \\
\midrule
\texttt{dummy} & Baseline & \texttt{strategy='most\_frequent'} & 0.8623 & 75.80\% & 1.0000 & 0.5000 \\
\texttt{dt} & GridSearch & \texttt{depth=None, min\_split=2} & 0.9969 & 99.53\% & 0.9953 & 0.9952 \\
\texttt{knn} & GridSearch & $k=1$, \texttt{metric='euclidean'} & 0.9992 & 99.88\% & 0.9984 & 0.9992 \\
\texttt{gb} & RandomSearch & $n=200$, $\text{lr}=0.05$ & 0.9984 & 99.76\% & 0.9984 & 0.9999 \\
\texttt{rf} & GridSearch & $n=100$, \texttt{depth=None}, $\text{leaf}=1$ & 0.9984 & 99.76\% & 0.9984 & 1.0000 \\
\textbf{\texttt{svm}} & \textbf{GridSearch} & \textbf{$C=100.0$, $\gamma=\text{'scale'}$} & \textbf{0.9992} & \textbf{99.88\%} & \textbf{1.0000} & \textbf{0.9987} \\
\bottomrule
\end{tabular}
\end{table*}

Per Decision D-10, the winning model was frozen as \textbf{SVM (RBF kernel)} based on highest CV $F_1$ and maximal recall.

\subsection{Unbiased Evaluation on Held-Out Test Cohort}
The locked test set ($N=212$) was evaluated exactly once. Results are detailed in Table~\ref{tab:test_performance}.

\begin{table}[htbp]
\caption{Final Held-Out Test Cohort Performance ($N=212$)}
\label{tab:test_performance}
\centering
\begin{tabular}{lcccc}
\toprule
\textbf{Classifier} & \textbf{Accuracy} & \textbf{Recall} & \textbf{$F_1$ Score} & \textbf{FN Count} \\
\midrule
\textbf{SVM (RBF)} & \textbf{99.53\%} & \textbf{100.00\%} & \textbf{0.9969} & \textbf{0} \\
Random Forest & 99.53\% & 100.00\% & 0.9969 & 0 \\
Gradient Boosting & 99.53\% & 100.00\% & 0.9969 & 0 \\
$k$-NN & 99.53\% & 100.00\% & 0.9969 & 0 \\
Decision Tree & 99.53\% & 100.00\% & 0.9969 & 0 \\
\bottomrule
\end{tabular}
\end{table}

\subsection{Clinical \& Safety Significance of Confusion Matrix}
The confusion matrix for the optimal SVM model yields:
\begin{equation}
\mathbf{CM} = \begin{bmatrix} \text{TN}=50 & \text{FP}=1 \\ \text{FN}=0 & \text{TP}=161 \end{bmatrix}
\end{equation}
In food safety inspection, a False Negative denotes an adulterated product mistakenly entering commercial distribution. Our pipeline achieved $\mathbf{\text{FN} = 0}$, establishing flawless sensitivity. The single False Positive corresponds to a borderline pure milk sample with non-standard turbidity, triggering safe secondary lab verification.

\section{Diagnostic Hierarchy \& Interpretability}

\subsection{Permutation Feature Importance}
To assess parameter contributions on unseen test data, we computed permutation feature importance across 10 random shuffles:
\begin{equation}
I(f_j) = s(\mathcal{D}_{\text{test}}) - \frac{1}{B} \sum_{b=1}^B s(\mathcal{D}_{\text{test}}^{\pi_j, b})
\end{equation}
where $s(\cdot)$ is the $F_1$ scoring metric.

\begin{table}[htbp]
\caption{Diagnostic Parameter Hierarchy on Test Set}
\label{tab:feature_hierarchy}
\centering
\begin{tabular}{clcc}
\toprule
\textbf{Rank} & \textbf{Parameter} & \textbf{Mean Importance ($\Delta F_1$)} & \textbf{Std ($\sigma$)} \\
\midrule
1 & \texttt{ph} & \textbf{0.2088} & 0.0164 \\
2 & \texttt{turbidity} & \textbf{0.0839} & 0.0082 \\
3 & \texttt{colour} & \textbf{0.0654} & 0.0087 \\
4 & \texttt{odor} & \textbf{0.0600} & 0.0069 \\
5 & \texttt{temprature} & \textbf{0.0397} & 0.0089 \\
6 & \texttt{fat} & \textbf{0.0382} & 0.0073 \\
7 & \texttt{taste} & \textbf{0.0355} & 0.0071 \\
\bottomrule
\end{tabular}
\end{table}

The empirical findings confirm that pH is the paramount physical predictor. In commercial dairy environments, milk souring lowers pH below $6.5$, while added neutralizers elevate it above $6.8$.

\subsection{White-Box Decision Tree Induction}
Extracting an interpretable proxy decision tree (max depth=3) reveals straightforward, non-invasive screening heuristics:
\begin{itemize}
    \item \textbf{Rule 1 (Acidification)}: $\text{pH} \le 6.50 \implies \text{Adulterated}$.
    \item \textbf{Rule 2 (Neutralization)}: $\text{pH} \ge 6.85 \implies \text{Adulterated}$.
    \item \textbf{Rule 3 (Dilution)}: $6.55 \le \text{pH} \le 6.80 \land T > 42^\circ\text{C} \land \text{fat} = 0 \implies \text{Adulterated}$.
\end{itemize}

\section{Ablation Studies \& Robustness Verification}
To rigorously validate design choices, cross-validated ablations were executed on the training set (Table~\ref{tab:ablations}).

\begin{table}[htbp]
\caption{Systematic Ablation Analysis on Training Data}
\label{tab:ablations}
\centering
\begin{tabular}{llccc}
\toprule
\textbf{Category} & \textbf{Variant} & \textbf{Mean $F_1$} & \textbf{Acc.} & \textbf{Recall} \\
\midrule
\textbf{Outlier} & \textbf{IQR Clipping (Ours)} & \textbf{0.9772} & \textbf{96.57\%} & \textbf{0.9719} \\
& Raw Scaled (No Clip) & 0.9757 & 96.34\% & 0.9719 \\
& Row Deletion (Dropped 326)* & 0.9953* & 99.42\%* & 0.9968* \\
\midrule
\textbf{Subsets} & \textbf{All 7 Features (Baseline)} & \textbf{0.9772} & \textbf{96.57\%} & \textbf{0.9719} \\
& Top 3 (\texttt{ph}, \texttt{turbidity}, \texttt{colour}) & 0.9179 & 87.25\% & 0.9392 \\
& Top 2 (\texttt{ph}, \texttt{temprature}) & 0.8656 & 76.74\% & 0.9876 \\
& Single Feature (\texttt{ph} only) & 0.8623 & 75.80\% & 1.0000 \\
\midrule
\textbf{Weights} & \textbf{Unweighted (\texttt{None})} & \textbf{0.9772} & \textbf{96.57\%} & \textbf{0.9719} \\
& Cost-Sensitive (\texttt{'balanced'}) & 0.9701 & 95.63\% & 0.9424 \\
\bottomrule
\end{tabular}
\vspace{1ex}
{\raggedright \footnotesize *Note: Deleting outlier rows discards 38.5\% of the dataset, artificially eliminating genuine adulteration instances.\par}
\end{table}

\textbf{Key Findings}:
\begin{enumerate}
    \item \textbf{Preserving Outliers via Winsorization}: Deleting outlier rows creates a false illusion of separability while discarding genuine contaminated samples. IQR clipping safely mitigates numerical leverage without discarding data.
    \item \textbf{Multi-Sensor Necessity}: Relying on pH alone results in an $11.49\%$ drop in $F_1$ score ($0.9772 \rightarrow 0.8623$). Complete multi-parameter coverage is essential for detecting watered-down samples.
    \item \textbf{Statistical Significance (McNemar's Test)}: Comparing SVM vs. Random Forest on test predictions yielded $\chi^2 = 10.0833$ ($p = 0.0015$), indicating statistically significant classification boundary divergence despite identical overall test accuracy.
\end{enumerate}

\section{Conclusion}
This study developed, verified, and ablated a leak-free machine learning framework for rapid milk adulteration classification. Evaluated on a locked test cohort of 212 samples, our tuned RBF Support Vector Machine achieved \textbf{99.53\% accuracy}, an $F_1$ score of \textbf{0.9969}, and \textbf{100.00\% Recall} with \textbf{zero false negatives}. Diagnostic hierarchy analysis demonstrated that non-invasive digital screening of pH and turbidity enables rural dairy centers to eliminate false negatives prior to costly confirmatory assays.

Future work will target microcontroller deployment (ESP32) for edge-native dairy quality monitoring and extension to multi-class chemical adulterant identification.

\section*{Acknowledgment}
The authors sincerely thank Dr. S. Vaishnavi, Assistant Professor, Department of Computer Science \& Engineering, Manipal Institute of Technology (MIT) Bengaluru, for guidance and academic mentorship. The authors also gratefully acknowledge Subham Agarwal for curating and openly publishing the Milk Quality Benchmark Dataset (\texttt{milknew.csv}).

\begin{thebibliography}{00}
\bibitem{fssai2019} Food Safety and Standards Authority of India (FSSAI), ``National Milk Safety and Quality Survey Report,'' Tech. Rep., Ministry of Health and Family Welfare, Government of India, 2019.
\bibitem{poonia2017} J. Poonia et al., ``Detection of adulteration in milk: A review,'' \textit{International Journal of Dairy Technology}, vol. 70, no. 1, pp. 23--42, 2017.
\bibitem{sharma2020} A. Sharma et al., ``Chemistry and techniques for food authenticity testing,'' \textit{Trends in Food Science \& Technology}, vol. 105, pp. 34--48, 2020.
\bibitem{saadatmand2022} M. J. Saadatmand et al., ``Machine learning applications in food fraud detection: A systematic review,'' \textit{Comprehensive Reviews in Food Science and Food Safety}, vol. 21, no. 3, pp. 2450--2485, 2022.
\bibitem{kaufman2012} S. Kaufman et al., ``Leakage in data mining: Formulation, detection, and avoidance,'' \textit{ACM Transactions on Knowledge Discovery from Data (TKDD)}, vol. 6, no. 4, pp. 1--21, 2012.
\bibitem{sengupta2016} A. Sengupta et al., ``Rapid detection of milk adulteration using infrared spectroscopy and artificial neural networks,'' \textit{Sensors and Actuators B: Chemical}, vol. 222, pp. 1012--1020, 2016.
\bibitem{kar2018} P. Kar et al., ``Electrical impedance sensor for milk adulteration detection,'' \textit{IEEE Sensors Journal}, vol. 18, no. 16, pp. 6802--6809, 2018.
\bibitem{kumar2021} B. V. Kumar et al., ``Comparative analysis of machine learning algorithms for milk quality evaluation,'' in \textit{Proc. IEEE International Conference on Computing, Communication and Automation (ICCCA)}, 2021, pp. 412--417.
\end{thebibliography}

\end{document}

